Shelfmark 103, scan 4.

\scanpage{4}\mspage{3}
\[
k(C)\handwritten{\in}H^{\handwritten{3}}(\handwritten{\pi}_{0}(C),\handwritten{\unsure{\pi|K}}_{1}(C))\ .
\]
By making these reflections more precise, one finds that the classification, up to $\otimes$-equivalence, of (Gr)-categories, is carried out precisely in terms of\deleted{\rlap{x}\unsure{e}} a group $\handwritten{\pi}_{0}$, of a commutative group $\handwritten{\pi}_{1}$ on which $\handwritten{\pi}_{0}$ acts, and of a cohomology class (which may be taken to be arbitrary) in $H^{\handwritten{3}}(\handwritten{\pi}_{0},\handwritten{\pi}_{1})$. The \deleted{\rlap{xxxxx}comp} group law on $H^{\handwritten{3}}$ moreover admits a ``geometric'' interpretation in the manner of Baer, in terms of operations on \handwritten{(}Gr)-categories.

\uline{Special cases}\quad In example 1, one \inserted{must} of course recover the first Postnikoff invariant $\handwritten{\in}H^{\handwritten{3}}(\handwritten{\pi}_{1}(X),\handwritten{\pi}_{2}(X))$, with a new ``geo\deleted{\unsure{(}}metric'' interpretation. In example 2, one finds $\handwritten{\pi}_{0}=H^{1}(X,F)$, $\handwritten{\pi}_{1}=H^{0}(X,F)$\handwritten{,} and the element of $H^{\handwritten{3}}$ is zero (because one has a ``strict Picard category'', cf\inserted{\handwritten{ Deligne}} \handwritten{[\, ]}). In \deleted{\rlap{xxx}\unsure{let|the}} the last example, one has $\handwritten{\pi}_{0}={}$ group of isomorphism classes of autoequivalences of $A$, $\handwritten{\pi}_{1}={}$ group of automorphisms of the identity functor of $A$, and the invariant $k(C)$ seems new (even in the case where $C$ is the category of invertible biModules over a Ring $\underline{O}_{X}$, where $\handwritten{\pi}_{0}$ can thus be interpreted as a ``Brandt groupoid'' and $\handwritten{\pi}_{1}$ as the group of ``central units'' of $\underline{O}_{X}$). Let us point out however that when $A$ is the category of torsors under an (ordinary) group $F$, one finds \deleted{\rlap{xxxxxxxxxxxxxxxxxxx}an element of $H^{\handwritten{3}}($} $\handwritten{\pi}_{0}={}$ group of outer automorphisms of $F$, $\handwritten{\pi}_{1}={}$ centre of $F$, and the element $k(C)$ of $H^{\handwritten{3}}$ is none other than the classical Mac-Lane invariant. It would remain to \deleted{\rlap{xxxxxxxx}\unsure{establish|establ}}\inserted{develop} a \inserted{geometric} \uline{direct} link between the theory of\deleted{\unsure{s|x}} Mac-Lane of extensions of noncommutative groups\deleted{\unsure{s|,}} and of bands (``kernels''\deleted{\unsure{)|,}} in his terminology) with which they are associated, and the theory \deleted{\rlap{xxxxxxxxxxxxxxxx}developed\unsure{s here,| here,}} sketched here of (Gr)-categ\deleted{\unsure{or|rr}}ories \unsure{at|and} of their splittings.

2. \uline{Picard categories}.

By a ``Picard category'' one means an ACU $\otimes$-category (associative-commutative-unital , cf. \handwritten{[\, ]}) which is a (\deleted{\unsure{E}}Gr)-category for its \deleted{\rlap{xxxx}\unsure{(Gr)|str.}} AU $\otimes$-structure. One says that it is a \uline{strict} Picard category if for every $X\handwritten{\in}\mathrm{Ob}\ C$, the symmetry of $X\otimes X$ is the identity. These \deleted{\rlap{xxxxx}categ} latter categories are studied systematically (in the more general context of ``Picard stacks'') by Deligne \handwritten{[\, ]}.

Let $C$ be a Picard category. As a (Gr)-category, $C$ possesses invariants $\handwritten{\pi}_{0}$ and $\handwritten{\pi}_{1}$, but in the present case $\handwritten{\pi}_{0}$ is also commutative and acts trivially on $\handwritten{\pi}_{\unsure{0|1}}$. One should note here that an equivalence of (Gr)-categories between two Picard categories $C,C'$ is not necessarily compatible with the commutativity constraints, thus is not necessarily an equivalence of categories of
% archive sheet number: 3
